% chapters/06/contoh.tex
\section{Dua Contoh Terapan}
\label{sec:contoh-bab-6}

\subsection{Contoh Pertama: Melacak Jatuh Bebas Bingkai demi Bingkai}

\begin{contohbox}
Sebuah benda dilepaskan tanpa kecepatan awal dengan gravitasi 1800 piksel per
detik kuadrat pada laju 60 bingkai per detik. Karena percepatan gravitasi
dikalikan selang waktu bernilai

\[
1800 \times \frac{1}{60} = 30 \text{ piksel per detik},
\]

kecepatan benda bertambah 30 piksel per detik pada setiap bingkai. Perpindahan
pada setiap bingkai adalah kecepatan pada bingkai itu dikalikan selang waktu.
\end{contohbox}

\begin{table}[htbp]
\centering
\caption{Posisi benda yang jatuh bebas, bingkai demi bingkai}
\label{tab:jatuh-bebas}
{\setlength{\tabcolsep}{4pt}\footnotesize%
\begin{tabularx}{\textwidth}{@{}C{1.6cm} C{3.0cm} C{3.4cm} Y@{}}
\toprule
\headrow \textbf{Bingkai} & \textbf{Kecepatan (piksel/detik)} & \textbf{Perpindahan bingkai itu} & \textbf{Posisi kumulatif (piksel)} \\
\midrule
1 & 30 & 0,5 piksel & 0,5 \\
\addlinespace
2 & 60 & 1,0 piksel & 1,5 \\
\addlinespace
3 & 90 & 1,5 piksel & 3,0 \\
\addlinespace
4 & 120 & 2,0 piksel & 5,0 \\
\addlinespace
5 & 150 & 2,5 piksel & 7,5 \\
\addlinespace
6 & 180 & 3,0 piksel & 10,5 \\
\bottomrule
\end{tabularx}}
\tabsource{Pada kolom ketiga, perpindahan bertambah 0,5 piksel setiap bingkai;
itulah wujud percepatan pada perhitungan yang terputus.}
\end{table}

Perhatikan pola pada kolom ketiga: perpindahannya bertambah 0,5 piksel setiap
bingkai. Itulah cara percepatan bekerja pada simulasi yang terputus: bukan
dengan mengubah posisi secara langsung, melainkan dengan menambah kecepatan
yang kemudian mengubah posisi.

Satu hal yang perlu dicatat sebagai perbandingan: menurut perhitungan yang
berlaku pada gerak bersinambung, benda itu seharusnya baru jatuh sejauh
$\frac{1}{2} \times 1800 \times (0{,}1)^2 = 9{,}0$ piksel setelah enam bingkai.
Simulasi menghasilkan 10,5 piksel, yaitu sedikit lebih cepat. Selisih itu
adalah galat integrasi Euler yang telah dibahas pada Bagian~2.4 dan diukur
pada Bagian~\ref{sec:program-fisika}: simulasi cenderung bergerak sedikit lebih
cepat daripada lintasan sebenarnya.

\subsection{Contoh Kedua: Menyusutkan Tinggi Pantulan}

\begin{contohbox}
Sebuah bola menumbuk lantai dengan kecepatan 700 piksel per detik dan
koefisien restitusi 0,8. Kecepatan setelah setiap pantulan menjadi 0,8 kali
kecepatan sebelumnya, sedangkan tinggi yang dicapai menjadi
$0{,}8^2 = 0{,}64$ kali tinggi sebelumnya, sebab tinggi sebanding dengan
kuadrat kecepatan.
\end{contohbox}

\begin{table}[htbp]
\centering
\caption{Kecepatan dan tinggi bola pada setiap pantulan}
\label{tab:pantulan-bertahap}
{\setlength{\tabcolsep}{4pt}\footnotesize%
\begin{tabularx}{\textwidth}{@{}C{2.2cm} C{3.0cm} C{2.8cm} Y@{}}
\toprule
\headrow \textbf{Pantulan} & \textbf{Kecepatan} & \textbf{Tinggi} & \textbf{Keterangan} \\
\midrule
Sebelum pertama & 700 piksel/detik & 136 piksel & Jatuh dari ketinggian ini. \\
\addlinespace
Pertama & 560 piksel/detik & 87 piksel & Kecepatan tinggal 0,8 kali. \\
\addlinespace
Kedua & 448 piksel/detik & 56 piksel & Tinggi tinggal 0,64 kali. \\
\addlinespace
Ketiga & 358 piksel/detik & 36 piksel & Tinggi tinggal seperempat tinggi semula. \\
\bottomrule
\end{tabularx}}
\end{table}

Setelah tiga pantulan, bola hanya mencapai ketinggian sekitar 36 piksel,
walaupun kecepatannya baru berkurang menjadi sekitar setengah. Perbandingan
itu menjelaskan mengapa pantulan pada permainan terasa cepat mereda: yang
menyusut adalah tinggi, dan tinggi menyusut secepat kuadrat kecepatan.

Apabila Anda menghendaki bola pantulan yang tetap terasa hidup lebih lama,
naikkan koefisien restitusinya, misalnya menjadi 0,9. Akan tetapi, ingatlah
bahwa benda dengan koefisien 1 tidak akan pernah berhenti, dan permainan
menjadi tidak pernah tenang.
